\documentclass[12pt,a4paper]{article}
\usepackage{cmap} % Makes the PDF copiable. See http://tex.stackexchange.com/a/64198/25761
\usepackage[T1]{fontenc}
\usepackage[brazil]{babel}
\usepackage[utf8]{inputenc}
\usepackage{amsmath}
\usepackage{amsfonts}
\usepackage{amssymb}
\usepackage{amsthm}
\usepackage{textcomp} % \degree
\usepackage{gensymb} % \degree
\usepackage[usenames,svgnames,dvipsnames]{xcolor}
\usepackage{hyperref}
\usepackage{graphicx}
\usepackage{systeme}
\usepackage[margin=2cm]{geometry}

\hypersetup{
    colorlinks = true,
    allcolors = {blue}
}

\newcommand{\vect}[1]{\overrightarrow{#1}}
\newcommand{\norm}[1]{\left|\left|{#1}\right|\right|}

\newcommand*\tipo{PROVA II}
\newcommand*\turma{TURMA E}
\newcommand*\disciplina{GAN0001}
\newcommand*\nome{GEOMETRIA ANALÍTICA}
\newcommand*\eu{Helder G. G. de Lima}
\newcommand*\data{30/04/2015}

\author{\eu}
\title{\tipo - \disciplina}
\date{\data}

\begin{document}
\thispagestyle{empty}
\newgeometry{margin=2cm,bottom=0.5cm}
\begin{center}
\includegraphics{udesc_joinville_cabecalho.pdf}
\\ Prof. \eu\footnote{
Este é um material de acesso livre distribuído sob os termos da licença \href{https://creativecommons.org/licenses/by-sa/4.0/deed.pt_BR}{Creative Commons BY-SA 4.0}.}

\noindent\begin{tabular}{l c c r}
  \textbf{\disciplina}
& \textbf{\tipo}
& \textbf{\data}
& \textbf{\turma}
\end{tabular}\vspace{-0.3cm}
\noindent\rule{17cm}{0.01cm}
\end{center}

\noindent Nome do aluno: \rule{14cm}{0.01cm}

\section*{Instruções}

\begin{enumerate}
\renewcommand{\theenumi}{\Roman{enumi}}
\item Identifique-se em todas as folhas.
\item Mantenha o celular e os demais equipamentos eletrônicos desligados durante a prova.
\item Justifique cada resposta com cálculos ou argumentos baseados na teoria estudada.
\item Escolha \textsc{\textbf{uma}} questão para \textsc{\textbf{não}} fazer (ela não será corrigida): \rule{3cm}{0.01cm}
\end{enumerate}

\section*{Questões}

\begin{enumerate}
\item (2,0 pontos) Dados $\vec{u}=(-1,3,0)$, $\vec{v}=(0,-4,1)$ e $\vec{w}=(1,0,1)$, calcule
\[
  \vec{u} \times (\vec{v} \times \vec{w}) \, + \,
  \vec{v} \times (\vec{w} \times \vec{u}) \, + \,
  \vec{w} \times (\vec{u} \times \vec{v}).
\]

\item (2,0 pontos) Encontre equações paramétricas da reta $r$ dada pela interseção do plano $-2x + 8y - 6z + 8 = 0$ com o plano $-x + 4y + 8z + 4 = 0$.

\item (2,0 pontos) Seja $s$ a reta $\frac{x - 2}{4} = \frac{y - 3}{-3} = z + 1$. Dê um exemplo (com equações paramétricas) de uma reta $r_1 // s$, de uma reta $r_2$ concorrente com $s$ e de uma reta $r_3$ reversa com $s$ (explique suas escolhas).


\item (2,0 pontos) Calcule o ângulo entre a reta $r : \begin{cases}
y = -x -1 \\
z = 2
\end{cases}$ e $s: x+1 = \dfrac{y-1}{-2} = z-1$.

\item (2,0 pontos) Forneça uma equação geral do plano que passa pelo ponto médio de $A = (-2,5,1)$ e $B = (4,-1,-7)$, e que além disso contém a reta $x = \frac{y}{-1} = z$.

\item (2,0 pontos) Determine as equações reduzidas (com variável independente $x$) da reta $r$ que passa por $A = (1,2,4)$ e que é perpendicular ao plano $\frac{x}{2} + y - 2z - 6 = 0$.
\end{enumerate}

\section*{Lembrete}
Se $\vec{u} = (a,b,c)$ então $\norm{\vec{u}} = \sqrt{a^2 + b^2 + c^2}$. Dados $\vec{v} = (a_1, b_1, c_1)$ e $\vec{w} = (a_2, b_2, c_2)$, tem-se:

\begin{itemize}
\item $\vec{v} \cdot \vec{w}
= a_1 a_2 + b_1 b_2 + c_1 c_2
= \norm{\vec{v}} \cdot \norm{\vec{w}} \cdot \cos(\theta)
$, onde $\theta$ é o ângulo entre os vetores $\vec{v}$ e $\vec{w}$.
\item $\vec{v} \times \vec{w}
= \begin{vmatrix}
\vec{i} & \vec{j} & \vec{k} \\
a_1 & b_1 & c_1\\
a_2 & b_2 & c_2
\end{vmatrix}$
\item $(\vec{u}, \vec{v}, \vec{w})
= \vec{u} \cdot \vec{v} \times \vec{w}
= \begin{vmatrix}
a   & b   & c \\
a_1 & b_1 & c_1 \\
a_2 & b_2 & c_2
\end{vmatrix}$
\end{itemize}

\newpage
\restoregeometry
\section*{Respostas e observações}

\begin{enumerate}
\item \textit{Dados $\vec{u}=(-1,3,0)$, $\vec{v}=(0,-4,1)$ e $\vec{w}=(1,0,1)$, calcule
\[
  \vec{u} \times (\vec{v} \times \vec{w}) \, + \,
  \vec{v} \times (\vec{w} \times \vec{u}) \, + \,
  \vec{w} \times (\vec{u} \times \vec{v}).
\]
}
Levando em conta que $\vec{u} \times (\vec{v} \times \vec{w})
= \begin{vmatrix}
\vec{v} & \vec{w} \\
\vec{u} \cdot \vec{v} & \vec{u} \cdot \vec{w}
\end{vmatrix}$, para quaisquer $\vec{u}, \vec{v}$ e $\vec{w}$, podemos calcular primeiramente:
\[
\vec{u} \cdot \vec{v} = (-1) \cdot 0 + 3 \cdot (-4) + 0 \cdot 1 = -12
\]
\[
\vec{u} \cdot \vec{w} = (-1) \cdot 1 + 3 \cdot 0 + 0 \cdot 1 = -1
\]
\[
\vec{v} \cdot \vec{w} = 0 \cdot 1 + (-4) \cdot 0 + 1 \cdot 1 = 1
\]
e então substituir:
\[
\vec{u} \times (\vec{v} \times \vec{w})
= \begin{vmatrix}
\vec{v} & \vec{w} \\
-12 & -1
\end{vmatrix}
= -\vec{v} +12 \vec{w}
\]
\[
\vec{v} \times (\vec{w} \times \vec{u})
= \begin{vmatrix}
\vec{w} & \vec{u} \\
1 & -12
\end{vmatrix}
= -12\vec{w} - \vec{u}
\]
\[
\vec{w} \times (\vec{u} \times \vec{v})
= \begin{vmatrix}
\vec{u} & \vec{v} \\
-1 & 1
\end{vmatrix}
= \vec{u} + \vec{v}
\]
Assim, a soma procurada é $(-\vec{v} + 12 \vec{w}) + (-12\vec{w} - \vec{u}) + (\vec{u} + \vec{v}) = \vec{0}$.

\textbf{Nota:} também é possível expandir expressões como $-12\vec{w} - \vec{u} = -12(1,0,1) - (-1,3,0) = (-12,0,-12) - (-1,3,0) = (-11,-3,-12)$. Ao somar os vetores obtidos nos três casos, o resultado também será zero. A diferença é que terá feito mais cálculos, sem necessidade. De fato, não é preciso calcular numericamente nem mesmo os produtos escalares, pois o mesmo raciocínio funciona se desenvolver a soma utilizando as expressões $\vec{u} \cdot \vec{v}$, $\vec{u} \cdot \vec{w}$ e $\vec{v} \cdot \vec{w}$ diretamente, sem troca-las por números (verifique!).


\item \textit{Encontre equações paramétricas da reta $r$ dada pela interseção do plano $-2x + 8y - 6z + 8 = 0$ com o plano $-x + 4y + 8z + 4 = 0$.}
Todos os pontos $(x, y, z) \in r$ devem satisfazer simultaneamente às equações dos planos, isto é,
\[
\systeme{
-2x + 8y - 6z + 8 = 0,
 -x + 4y + 8z + 4 = 0}
\]
Da segunda equação, resulta que $x = 4y + 8z + 4$, e essa expressão pode ser substituída na primeira equação para obter que $-2(4y + 8z + 4) + 8y - 6z + 8 = 0$. Simplificando, conclui-se que $- 22z = 0$, o que significa que $z = 0$.

Com isso, a segunda equação do sistema pode ser reescrita como $-x + 4y + 4 = 0$, o que implica que $x = 4y + 4$. Assim, a reta $r$ pode ser descrita pelas equações reduzidas (com variável independente $y$)
\[
\systeme*{
x = 4y + 4,
z = 0.
}
\]
ou equivalentemente, pelas equações paramétricas
\[
\systeme*{
x = 4t + 4,
y = t,
z = 0
}
\]
que são obtidas considerando o ponto $A = (4, 0, 0) \in r$ e o vetor diretor $\vec{v} = (4, 1, 0)$ de $r$.

\item \textit{ Seja $s$ a reta $\frac{x - 2}{4} = \frac{y - 3}{-3} = z + 1$. Dê um exemplo (com equações paramétricas) de uma reta $r_1 // s$, de uma reta $r_2$ concorrente com $s$ e de uma reta $r_3$ reversa com $s$ (explique suas escolhas). }

A reta $s$ tem vetor diretor $\vec{v} = (4,-3,1)$.
\begin{enumerate}

\item Para obter uma reta $r_1 // s$, basta utilizar o mesmo vetor diretor $\vec{v}$ (ou um múltiplo dele), e qualquer ponto que não pertença a $s$ (pois os vetores diretores de retas paralelas são colineares, e elas só têm algum ponto em comum se forem coincidentes). Se escolhermos, por exemplo, $A = (0,0,0)$, é claro que suas coordenadas não satisfazem as equações simétricas de $s$, e portanto $A \not\in s$. As equações paramétricas da reta $r_1$ que passa por $A$ e que tem a direção de $\vec{v}$ são:
\[
\systeme*{
x = 4t,
y = -3t,
z = t.
}
\]

\item Para qualquer ponto $B \in s$ e qualquer vetor $\vec{u}$ que não seja colinear com $\vec{v}$, a reta que passa por $B$ e que tem a direção de $\vec{u}$ é concorrente com $s$. Assim, como $B = (2,3,-1) \in s$ e $\vec{u} = (1,0,0)$ não é múltiplo de $\vec{v}$, as equações paramétricas
\[
\systeme*{
x = 2 + t,
y = 3,
z = -1
}
\]
representam uma reta $r_2$, entre as infinitas retas que são concorrentes com $s$.

\item Se no item anterior trocarmos o ponto $B \in s$ por um ponto $C$ que não pertence nem a $s$ nem a $r_2$, obteremos uma reta $r_3$ que não é paralela a $s$, e que também não possui um ponto em comum com $s$, ou seja, $r_3$ e $s$ serão retas reversas. A título de exemplo, pode-se considerar $C = A = (0,0,0)$, que não pertence a $s$ nem a $r_2$ (pois não satisfaz suas equações), e obter as equações paramétricas
\[
\systeme*{
x = t,
y = 0,
z = 0
}
\]
que representam o eixo $Ox$.
\end{enumerate}


\item \textit{ Calcule o ângulo entre a reta $r : \begin{cases}
y = -x -1 \\
z = 2
\end{cases}$ e $s: x+1 = \dfrac{y-1}{-2} = z-1$. }

O vetor $\vec{u} = (1,-2,2) - (0,-1,2) = (1,-1,0)$ é um vetor diretor da reta $r$, e o vetor $\vec{v} = (1,-2,1)$ é um vetor diretor de $s$. Então, o ângulo $\theta$ entre $r$ e $s$ é dado por
\[
\cos(\theta)
= \frac{|\vec{u} \cdot \vec{v}|}{\norm{u} \cdot \norm{v}}
= \frac{|1 \cdot 1 + (-1) \cdot (-2) + 0 \cdot 1|}
       {\sqrt{1 + 1 + 0} \cdot \sqrt{1+4+1}}
= \frac{|3|}{\sqrt{2} \cdot \sqrt{6}}
= \frac{3}{2\sqrt{3}}
= \frac{\sqrt{3}}{2}.
\]
Logo, $\theta = \operatorname{arccos}(\frac{\sqrt{3}}{2}) = \frac{\pi}{6}$, que corresponde a um ângulo de $30\degree$.

\item \textit{ Forneça uma equação geral do plano que passa pelo ponto médio de $A = (-2,5,1)$ e $B = (4,-1,-7)$, e que além disso contém a reta $x = \frac{y}{-1} = z$. }

O ponto médio de $A$ e $B$ é $M = \left(\frac{-2+4}{2}, \frac{5-1}{2}, \frac{1-7}{2}\right) = (1,2,-3)$. O vetor $\vec{n}$, normal do plano $\pi$ procurado, será ortogonal a um vetor diretor $\vec{v} = (1,-1,1)$ da reta. Além disso, $\vec{n}$ também será ortogonal a qualquer vetor $\vec{u}$ que tenha origem em um ponto arbitrário da reta (por exemplo, $(1,-1,1)$), e extremidade em $M$. Deste modo, podemos obter um vetor $\vec{n}$ fazendo o produto vetorial entre $\vec{v}$ e $\vec{u} = M - (1,-1,1) = (0,3,-4)$:

\[
\vec{v} \times \vec{u} =
\begin{vmatrix}
\vec{i} & \vec{j} & \vec{k} \\
1 & -1 & 1 \\
0 & 3 & -4
\end{vmatrix}
= (4 - 3) \vec{i} - (-4 - 0) \vec{j} + (3 - 0) \vec{k}
= (1, 4, 3).
\]
Logo, a equação geral do plano tem a forma $x + 4y + 3z + d = 0$, e como $M \in \pi$, devemos ter $1 + 4 \cdot 2 + 3 \cdot(-3) + d = 0$, o que implica que $d = 0$. Portanto o plano tem equação geral $x + 4y + 3z = 0$.

\item \textit{ Determine as equações reduzidas (com variável independente $x$) da reta $r$ que passa por $A = (1,2,4)$ e que é perpendicular ao plano $\frac{x}{2} + y - 2z - 6 = 0$. }

O plano em questão é ortogonal ao vetor $\vec{n} = (\frac{1}{2}, 1, -2)$ (e também ao seu múltiplo $\vec{n}_1 = 2 \vec{n} = (1, 2, -4)$, caso prefira evitar frações). Utilizando o ponto $A$ e o vetor $\vec{n}$, conclui-se que a reta $r$ tem as seguintes equações paramétricas:
\[
\begin{cases}
x = 1 + \frac{t}{2}\\
y = 2 + t\\
z = 4 - 2t
\end{cases}
\]
Da primeira equação, resulta que $t = 2x - 2$, e então as demais equações podem ser reescritas assim:
\[
\begin{cases}
y = 2 + (2x-2)\\
z = 4 - 2(2x-2).
\end{cases}
\]
Portanto, as equações reduzidas são
\[
\begin{cases}
y = 2x\\
z = -4x + 8.
\end{cases}
\]

\end{enumerate}

\end{document}
